\documentclass{article}

\usepackage{graphicx} % Required for inserting images
\usepackage{xcolor}

\usepackage{amsmath}
\usepackage{amsthm}
\usepackage{thmtools} 
\usepackage{cleveref}


% Number equations within sections (instead of chapters)
\numberwithin{equation}{section}

% Declare theorem environments anchored to sections
\declaretheorem[name=Theorem,numberwithin=section]{theorem}
\declaretheorem[name=Lemma,numberlike=theorem]{lemma}
\declaretheorem[name=Proposition,numberlike=theorem]{proposition}
\declaretheorem[name=Corollary,numberlike=theorem]{corollary}
\declaretheorem[name=Definition,numberlike=theorem]{definition}
\declaretheorem[name=Assumption,numberlike=theorem]{assumption}
\declaretheorem[name=Example,numberlike=theorem]{example}
\declaretheorem[name=Remark,numberlike=theorem]{remark}

% For restatable theorems 
\declaretheorem[name=Theorem,numberwithin=section]{restatabletheorem}

\title{Robust Program Equilibrium}
\author{Colomban Duclaux}
\date{\today}

\usepackage[backend=biber,style=alphabetic]{biblatex}
\addbibresource{../references.bib}

\begin{document}

\maketitle

These notes recap the contents exposed in \cite{robust_program_equilibrium}
\section{Abstract}
Tennenholtz's program equilibrium is too fragile, and unrealistic. We propose another program:  cooperate with some small probability $\epsilon$ and otherwise act as the opponent
acts against this program. This is similar to the tit-for-tat strategy for the iterated prisoner's dilemma.
We generalize this one-shot versions to other games.

\section{Introduction}
Tennenholtz introduced the concepts of programm equilibria and proposed a framework that was too fragile.
We then move on to more robust program equilibria by letting the programs reason about each other, like \cite{critch2022cooperativeuncooperativeinstitutiondesigns}, but they are computationally expensive.
This paper proposes a more computationally feasable program equilibria, similar to \texttt{FairBot}.

\section{Robust program equilibrium in the prisoner's dilemma}
$\epsilon \texttt{GroundedFairBot}$: cooperate with a probability $\epsilon$, otherwise act as your opponent plays against you.
$2$ ideas: \begin{itemize}
    \item It is a version of $\texttt{FairBot}$ which mimics a tit-for-tat strategy
    \item It cooperates with a small probability to avoid running into the infinite loop that a $\texttt{FairBot}$ might encounter.
\end{itemize}

\section{Extensions}
$\epsilon \texttt{Grounded} \pi_i\texttt{Bot}$ $\to$ Does terminate.

Relationship to the underlying iterated game strategy.

\printbibliography

\end{document}